%%=============================================================================
%% CAPITULO 1 -- INTRODUCAO
%%
%% Na introducao, faca a apresentacao do trabalho. Ela devera ocupar uma ou,
%% no maximo, duas paginas. De um panorama da tematica, exponha a questao
%% problema que conduzira a pesquisa proposta, a hipotese e a justificativa.
%%
%% Titulos: nao utilize ponto apos a numeracao das secoes.
%%=============================================================================

\section{Introdução}

Na primeira página da parte textual do trabalho de conclusão de curso, faça a
apresentação de sua pesquisa. A introdução deverá ocupar uma ou, no máximo, duas
páginas do seu trabalho. Nela, dê um panorama da temática, exponha a questão
problema que conduzirá a pesquisa proposta, a hipótese e a justificativa.

Os capítulos (seções) e subcapítulos (subseções) são sequenciais. O título de
cada seção primária (\textbf{1 INTRODUÇÃO}, \textbf{2 REFERENCIAL TEÓRICO} etc.)
é grafado em letras maiúsculas e negritadas. Não utilize ponto após a numeração
das seções. Deixe uma linha em branco entre as seções.

\subsection{Objetivo geral}

Enuncie, em um único parágrafo e com um verbo no infinitivo, o que a pesquisa
pretende alcançar.

\subsection{Objetivos específicos}

\begin{itemize}
  \item Objetivo específico 1;
  \item Objetivo específico 2;
  \item Objetivo específico 3.
\end{itemize}

\subsection{Problema}

Formule a questão problema em forma de pergunta, de modo claro e delimitado.

\subsection{Hipótese}

Apresente a resposta provisória que a pesquisa pretende confirmar ou refutar.

\subsection{Justificativa}

Explique a relevância acadêmica e social do tema e o que motivou a escolha do
recorte adotado.
